% !PRINTED 624
\weberpage{624}
Die Gruppe $\Psi$ ist isomorph mit einer in $\mathfrak{N}$ (\S.\,169) enthaltenen Gruppe nach dem Modul $m$ genommener ganzer
rationaler Zahlen, die wir mit $\mathfrak{A}$ bezeichnen, die aus allen, einer
Congruenz
\weq{3}{a \equiv p^h \pmod{m}}
gen\"ugenden Zahlen $a$ besteht, und die wir daher symbolisch durch
\weq{4}{\mathfrak{A} \equiv p^h \pmod{m}}
darstellen k\"onnen, wenn $h$ einen beliebigen ganzzahligen Exponenten bedeutet (vergl. \S.\,18).

Wir k\"onnen also, wenn wir die Zahlen $\xi_1, \xi_2, \ldots, \xi_e$ aus $\mathfrak{N}$
passend ausw\"ahlen:
\weq{5}{\mathfrak{N} = \mathfrak{A}\xi_1 + \mathfrak{A}\xi_2 + \mathfrak{A}\xi_3 + \cdots + \mathfrak{A}\xi}
% [sic] source prints the final \xi without the expected subscript e; (6)--(7) supply the context
setzen, und jeder dieser Nebengruppen entspricht eines der conjugirten Functionale $\mathfrak{p}_{\xi_1}, \mathfrak{p}_{\xi_2}, \ldots, \mathfrak{p}_{\xi_e}$, die alle von einander verschieden sind und deren Product mit $p$ associirt ist.
% [sic] "Functionale" as printed; context elsewhere on the page uses "Primfactoren"/"Primideale" for this notion

Wir setzen also
\weq{6}{p = \mathfrak{p}_{\xi_1} \mathfrak{p}_{\xi_2} \ldots \mathfrak{p}_{\xi_e}.}
Um aber das Zahlensystem
\weq{7}{\xi_1, \xi_2, \ldots, \xi_e,}
auf das es haupts\"achlich ankommt, genauer zu charakterisiren,
bemerken wir, dass wir dazu jedes System von $e$ Zahlen der
Gruppe $\mathfrak{N}$ nehmen k\"onnen, wenn nur keine zwei unter ihnen,
$\xi, \xi'$, eine Congruenz
\[
\xi'\xi^{-1} \equiv p^h \pmod{m}
\]
erf\"ullen.

Um ein solches Zahlensystem zu finden, ist es gut, den Fall
eines ungeraden $m$ von dem anderen zu trennen, in dem $m$ eine
Potenz von 2 ist.

1. Ist zun\"achst $m = q^{\varkappa}$ ungerade, so nehmen wir eine
primitive Wurzel $g$ von $m$ (\S.\,14) und setzen
\weq{8}{p \equiv g^{\gamma} \pmod{m}.}
Denn $f$ ist [nach (1), (2)] die kleinste positive Zahl, die der
Bedingung
\[
\gamma f \equiv 0 \pmod{\mu}
\]
gen\"ugt, und folglich ist $e$ der gr\"osste gemeinschaftliche Theiler
von $\mu$ und $\gamma$, und irgend eine Zahl $n \equiv g^{\lambda}$ ist nur dann mit
%CONT>>
